% TOP_PROBLEM: {"id": 30006864, "title": "Worst-case to average-case hardness reduction for the Learning Parity with Noise problem", "category_id": 15, "category": {"id": 15, "name": "computer_science", "display_name": "Computer Science", "description": "Computational complexity, algorithms, and theoretical CS.", "slug": "computer-science", "order_index": 15, "created_at": "2026-07-31T15:26:25.671Z"}, "status": "open"}
% FIRST_POSTED: 2026-09-25
% ENTRY_KIND: statement_only
% !TeX program = lualatex
% Definition and sources only; this file is not an LLM solution attempt.
\documentclass[11pt]{article}
\usepackage[margin=25mm]{geometry}
\usepackage{fontspec}
\setmainfont{Latin Modern Roman}
\usepackage{amsmath,amssymb,mathtools}
\usepackage{unicode-math}
\setmathfont{Latin Modern Math}
\usepackage{xurl,hyperref}
\hypersetup{colorlinks=true,urlcolor=blue}
\setcounter{secnumdepth}{0}
\setlength{\parindent}{0pt}
\setlength{\parskip}{5pt}
\setlength{\emergencystretch}{3em}
\begin{document}
\section*{\texorpdfstring{Worst-case to average-case hardness reduction for the Learning Parity with Noise problem}{Worst-case to average-case hardness reduction for the Learning Parity with Noise problem}}\label{30006864}
\subsection{Definitions and mathematical statement}
For a secret $s\in{\mathbb{F}}_2^n$ and a fixed noise rate $0<\eta<1/2$, a Learning Parity with Noise sample is
\[
 (a,b),\qquad a\sim\operatorname{Unif}({\mathbb{F}}_2^n),\quad
 b=\langle a,s\rangle+e\pmod2,\quad e\sim\operatorname{Bernoulli}(\eta),
\]
with independent samples and errors. Search LPN recovers $s$ with significant probability from polynomially many samples; decision LPN distinguishes these samples from uniform pairs. The task asks for a worst-case-to-average-case hardness reduction from a worst-case decoding problem to standard constant-noise LPN. A decoding input typically specifies a binary code and a received word, with a distance promise, and asks for a nearby codeword. The catalogue does not select the source promise, approximation factor, reduction type, or search/decision version. These choices must be fixed in any claimed reduction; a low-noise result alone does not provide the stated constant-noise conclusion.
\par\smallskip\textit{Formulation and concept references:} \href{https://eprint.iacr.org/2026/1168}{[S1]}.

\subsection{Short English statement}
Can hardness of decoding arbitrary worst-case code instances be transferred rigorously to learning random noisy parities with a fixed positive noise rate?
\subsection{Sources}
\begin{itemize}
\item[S1]Aggarwal, D., Gupta, R., Nguyen, H.H., Tan, K.Z., Vasudevan, P.N. (2026). Towards Worst-case Hardness for Low-Noise LPN. IACR Cryptology ePrint Archive, Report 2026/1168.
\url{https://eprint.iacr.org/2026/1168}.
\textit{Formulation source.}
\end{itemize}
\end{document}
